% Options for packages loaded elsewhere
\PassOptionsToPackage{unicode}{hyperref}
\PassOptionsToPackage{hyphens}{url}
\documentclass[
]{article}
\usepackage{xcolor}
\usepackage{amsmath,amssymb}
\setcounter{secnumdepth}{-\maxdimen} % remove section numbering
\usepackage{iftex}
\ifPDFTeX
  \usepackage[T1]{fontenc}
  \usepackage[utf8]{inputenc}
  \usepackage{textcomp} % provide euro and other symbols
\else % if luatex or xetex
  \usepackage{unicode-math} % this also loads fontspec
  \defaultfontfeatures{Scale=MatchLowercase}
  \defaultfontfeatures[\rmfamily]{Ligatures=TeX,Scale=1}
\fi
\usepackage{lmodern}
\ifPDFTeX\else
  % xetex/luatex font selection
\fi
% Use upquote if available, for straight quotes in verbatim environments
\IfFileExists{upquote.sty}{\usepackage{upquote}}{}
\IfFileExists{microtype.sty}{% use microtype if available
  \usepackage[]{microtype}
  \UseMicrotypeSet[protrusion]{basicmath} % disable protrusion for tt fonts
}{}
\makeatletter
\@ifundefined{KOMAClassName}{% if non-KOMA class
  \IfFileExists{parskip.sty}{%
    \usepackage{parskip}
  }{% else
    \setlength{\parindent}{0pt}
    \setlength{\parskip}{6pt plus 2pt minus 1pt}}
}{% if KOMA class
  \KOMAoptions{parskip=half}}
\makeatother
\usepackage{color}
\usepackage{fancyvrb}
\newcommand{\VerbBar}{|}
\newcommand{\VERB}{\Verb[commandchars=\\\{\}]}
\DefineVerbatimEnvironment{Highlighting}{Verbatim}{commandchars=\\\{\}}
% Add ',fontsize=\small' for more characters per line
\newenvironment{Shaded}{}{}
\newcommand{\AlertTok}[1]{\textcolor[rgb]{1.00,0.00,0.00}{\textbf{#1}}}
\newcommand{\AnnotationTok}[1]{\textcolor[rgb]{0.38,0.63,0.69}{\textbf{\textit{#1}}}}
\newcommand{\AttributeTok}[1]{\textcolor[rgb]{0.49,0.56,0.16}{#1}}
\newcommand{\BaseNTok}[1]{\textcolor[rgb]{0.25,0.63,0.44}{#1}}
\newcommand{\BuiltInTok}[1]{\textcolor[rgb]{0.00,0.50,0.00}{#1}}
\newcommand{\CharTok}[1]{\textcolor[rgb]{0.25,0.44,0.63}{#1}}
\newcommand{\CommentTok}[1]{\textcolor[rgb]{0.38,0.63,0.69}{\textit{#1}}}
\newcommand{\CommentVarTok}[1]{\textcolor[rgb]{0.38,0.63,0.69}{\textbf{\textit{#1}}}}
\newcommand{\ConstantTok}[1]{\textcolor[rgb]{0.53,0.00,0.00}{#1}}
\newcommand{\ControlFlowTok}[1]{\textcolor[rgb]{0.00,0.44,0.13}{\textbf{#1}}}
\newcommand{\DataTypeTok}[1]{\textcolor[rgb]{0.56,0.13,0.00}{#1}}
\newcommand{\DecValTok}[1]{\textcolor[rgb]{0.25,0.63,0.44}{#1}}
\newcommand{\DocumentationTok}[1]{\textcolor[rgb]{0.73,0.13,0.13}{\textit{#1}}}
\newcommand{\ErrorTok}[1]{\textcolor[rgb]{1.00,0.00,0.00}{\textbf{#1}}}
\newcommand{\ExtensionTok}[1]{#1}
\newcommand{\FloatTok}[1]{\textcolor[rgb]{0.25,0.63,0.44}{#1}}
\newcommand{\FunctionTok}[1]{\textcolor[rgb]{0.02,0.16,0.49}{#1}}
\newcommand{\ImportTok}[1]{\textcolor[rgb]{0.00,0.50,0.00}{\textbf{#1}}}
\newcommand{\InformationTok}[1]{\textcolor[rgb]{0.38,0.63,0.69}{\textbf{\textit{#1}}}}
\newcommand{\KeywordTok}[1]{\textcolor[rgb]{0.00,0.44,0.13}{\textbf{#1}}}
\newcommand{\NormalTok}[1]{#1}
\newcommand{\OperatorTok}[1]{\textcolor[rgb]{0.40,0.40,0.40}{#1}}
\newcommand{\OtherTok}[1]{\textcolor[rgb]{0.00,0.44,0.13}{#1}}
\newcommand{\PreprocessorTok}[1]{\textcolor[rgb]{0.74,0.48,0.00}{#1}}
\newcommand{\RegionMarkerTok}[1]{#1}
\newcommand{\SpecialCharTok}[1]{\textcolor[rgb]{0.25,0.44,0.63}{#1}}
\newcommand{\SpecialStringTok}[1]{\textcolor[rgb]{0.73,0.40,0.53}{#1}}
\newcommand{\StringTok}[1]{\textcolor[rgb]{0.25,0.44,0.63}{#1}}
\newcommand{\VariableTok}[1]{\textcolor[rgb]{0.10,0.09,0.49}{#1}}
\newcommand{\VerbatimStringTok}[1]{\textcolor[rgb]{0.25,0.44,0.63}{#1}}
\newcommand{\WarningTok}[1]{\textcolor[rgb]{0.38,0.63,0.69}{\textbf{\textit{#1}}}}
\setlength{\emergencystretch}{3em} % prevent overfull lines
\providecommand{\tightlist}{%
  \setlength{\itemsep}{0pt}\setlength{\parskip}{0pt}}
\usepackage{microtype}
\usepackage{titlesec}
\usepackage{booktabs}
\usepackage{array}
\usepackage{longtable}
\usepackage{ragged2e}
\usepackage{etoolbox}
\usepackage{caption}
\usepackage{fancyhdr}
\usepackage{setspace}
\usepackage{newunicodechar}

\newunicodechar{✓}{$\checkmark$}
\newunicodechar{∈}{$\in$}
\newunicodechar{∃}{$\exists$}
\newunicodechar{≠}{$\neq$}
\newunicodechar{≤}{$\le$}
\newunicodechar{≥}{$\ge$}
\newunicodechar{→}{$\rightarrow$}
\newunicodechar{₁}{$_1$}
\newunicodechar{₂}{$_2$}
\newunicodechar{ₖ}{$_k$}
\newunicodechar{─}{-}
\newunicodechar{└}{+}
\newunicodechar{├}{+}
\newunicodechar{│}{|}

\setstretch{1.03}
\setlength{\parindent}{0pt}
\setlength{\parskip}{0.45em}
\setlength{\emergencystretch}{3em}

\titleformat{\section}{\Large\bfseries}{\thesection.}{0.6em}{}
\titleformat{\subsection}{\large\bfseries}{\thesubsection.}{0.6em}{}
\titleformat{\subsubsection}{\normalsize\bfseries}{\thesubsubsection.}{0.6em}{}

\titlespacing*{\section}{0pt}{1.4em}{0.6em}
\titlespacing*{\subsection}{0pt}{1.1em}{0.4em}
\titlespacing*{\subsubsection}{0pt}{0.8em}{0.3em}

\setlength{\LTpre}{0.35em}
\setlength{\LTpost}{0.35em}
\setlength{\LTleft}{0pt}
\setlength{\LTright}{0pt}
\setlength{\tabcolsep}{4pt}
\renewcommand{\arraystretch}{1.08}
\AtBeginEnvironment{longtable}{\small}
\newcolumntype{L}[1]{>{\RaggedRight\arraybackslash}p{#1}}
\newcolumntype{C}[1]{>{\Centering\arraybackslash}p{#1}}
\newcolumntype{R}[1]{>{\RaggedLeft\arraybackslash}p{#1}}

\captionsetup[table]{skip=4pt}
\captionsetup[figure]{skip=4pt}

\pagestyle{fancy}
\fancyhf{}
\fancyfoot[C]{\thepage}
\renewcommand{\headrulewidth}{0pt}
\renewcommand{\footrulewidth}{0pt}
\setlength{\headheight}{14pt}

\newcommand{\codeid}[1]{{\small\begingroup\urlstyle{tt}\nolinkurl{#1}\endgroup}}
\usepackage{bookmark}
\IfFileExists{xurl.sty}{\usepackage{xurl}}{} % add URL line breaks if available
\urlstyle{same}
\hypersetup{
  hidelinks,
  pdfcreator={LaTeX via pandoc}}

\author{}
\date{}

\begin{document}

\section{Ontology-Grounded Retrieval for Auditable LLM-Assisted Secure
Code
Generation}\label{ontology-grounded-retrieval-for-auditable-llm-assisted-secure-code-generation}

\textbf{Pedro Farinha}\\
Independent Researcher\\
\href{mailto:pedro.farinha@shiftleft.pt}{\nolinkurl{pedro.farinha@shiftleft.pt}}\\
\href{https://github.com/pedrofarinhaatshiftleftpt}{github.com/pedrofarinhaatshiftleftpt}

\emph{Supported by Shiftleft - Secure Software Engineering, lda.}

\begin{center}\rule{0.5\linewidth}{0.5pt}\end{center}

\subsection{Abstract}\label{abstract}

Large language models (LLMs) generate code with security vulnerabilities
at significant rates {[}1, 2, 3{]}. Recent work shows that
security-focused prompting can reduce vulnerability rates {[}4{]},
suggesting that LLMs can produce more secure code \emph{when given
structured security context} --- but current retrieval approaches lack
the requirement-level grounding needed to do this systematically.

We propose \textbf{ontology-grounded retrieval}, a method for providing
LLMs with typed, weighted, and provenanced security knowledge drawn from
a normalized domain ontology. Unlike plain vector-similarity RAG ---
which retrieves relevant text without provenance, epistemic weight, or
completeness guarantees --- ontology-grounded retrieval operates through
structured queries against a typed knowledge graph, producing a
retrieval result with six formally distinguishable properties:
first-class provenance, explicit typing, epistemic weighting,
reproducible grounding, assessable completeness, and verifiability
against the same index.

We define a \textbf{retrieval contract} with completeness and provenance
invariants, present a \textbf{verification taxonomy} (syntactic,
semantic, human) that makes automation boundaries explicit, and
illustrate the method through three compact scenarios. The validation is
methodological, not experimental; a controlled evaluation is identified
as future work.

The method does not make LLM output secure. It makes the grounding
\textbf{auditable}: for any generated code, one can determine which
security requirements informed it, at what authority level, and whether
all applicable requirements were addressed.

\textbf{Keywords:} LLM, code generation, application security, ontology,
retrieval-augmented generation, structured retrieval, requirements
traceability, verification

\begin{center}\rule{0.5\linewidth}{0.5pt}\end{center}

\subsection{1. Introduction}\label{introduction}

LLM-assisted code generation has achieved wide adoption {[}5{]}, but the
security properties of generated code remain uncontrolled. Evaluations
across multiple models and languages show vulnerability rates of
25--40\% {[}1, 2, 3{]}, spanning injection, authentication,
cryptographic misuse, and configuration errors.

A natural mitigation is to inject security documentation into the LLM's
context via retrieval-augmented generation (RAG) {[}6{]}. If the LLM has
access to relevant security guidelines, it should produce more secure
code. In practice, however, this approach provides security context
through vector similarity search --- a mechanism that is probabilistic,
untyped, and provenance-free. The LLM receives text that is similar to
the query but has no structured understanding of:

\begin{itemize}
\tightlist
\item
  \textbf{Which requirement} the text represents
\item
  \textbf{At what authority level} --- mandatory requirement or
  informational example?
\item
  \textbf{Whether all relevant requirements} were retrieved ---
  similarity has no concept of domain coverage
\item
  \textbf{How to verify} that the output satisfies the retrieved
  requirements
\end{itemize}

This paper proposes \textbf{ontology-grounded retrieval} as a
methodological alternative. The method uses a normalized security
ontology as the retrieval substrate, producing typed, weighted,
provenanced context that enables a closed verification loop. The method
is general: it requires a typed ontology with explicit relations, a
stratified structural index, and a delivery protocol. We instantiate it
using AppSec Core {[}7{]} (a normalization ontology for application
security) and the Model Context Protocol (MCP) {[}8{]}, but the method
is not bound to these specific components.

\subsubsection{Research Question}\label{research-question}

\textbf{How can structured, ontology-based retrieval improve
reproducibility, completeness, and traceability in LLM-assisted secure
code generation --- compared to plain vector-similarity retrieval?}

\subsubsection{Contributions}\label{contributions}

\begin{enumerate}
\def\labelenumi{\arabic{enumi}.}
\item
  \textbf{Ontology-grounded retrieval as a method}: a five-stage
  pipeline (intent parsing → structured retrieval → context assembly →
  grounded generation → verification) with formally distinguishable
  properties.
\item
  \textbf{A retrieval contract} with completeness and provenance
  invariants that make grounding testable.
\item
  \textbf{A six-property comparison} between plain vector-similarity RAG
  and ontology-grounded retrieval (provenance, typing, weighting,
  reproducibility, completeness, verifiability).
\item
  \textbf{A three-class verification taxonomy} (syntactic, semantic
  approximation, human judgment) that makes automation boundaries
  explicit.
\item
  \textbf{Three compact illustrative scenarios} showing how the audit
  trail is intended to operate from prompt through verification.
\end{enumerate}

\subsubsection{Scope}\label{scope}

This paper describes a \textbf{method, its contract properties, and its
AppSec-specific instantiation}. We do not claim empirical superiority
over alternative retrieval architectures; we claim that the proposed
structure makes grounding auditable and its completeness properties
assessable. The method is general in structure but validated here only
in an AppSec instantiation using AppSec Core {[}7{]} and MCP {[}8{]}.
Generality beyond this domain is argued, not demonstrated. A controlled
evaluation is future work.

\begin{center}\rule{0.5\linewidth}{0.5pt}\end{center}

\subsection{2. Background and Related
Work}\label{background-and-related-work}

\subsubsection{2.1 LLM Code Generation and
Security}\label{llm-code-generation-and-security}

Modern LLMs generate syntactically and often semantically correct code
{[}2, 5{]}, but optimize for functional correctness, not security.
Pearce et al.~{[}1{]} found that \textasciitilde40\% of
Copilot-generated programs contained vulnerabilities. Tihanyi et
al.~{[}2{]} confirmed this at scale across 9 LLMs and 331,000 programs.
Tony et al.~{[}3{]} developed LLMSecEval as a standardized security
benchmark.

Tony et al.~{[}4{]} demonstrated that security-focused prompting
(recursive criticism, chain-of-thought with security examples) can
reduce vulnerability rates by up to 75\% on some tasks. This is a key
observation: \textbf{LLMs can produce more secure code when given
appropriate security context}. The question is how to provide that
context systematically, completely, and auditably.

\subsubsection{2.2 Retrieval-Augmented
Generation}\label{retrieval-augmented-generation}

Lewis et al.~{[}6{]} introduced RAG for grounding LLM outputs in
external knowledge. RAG is a family of architectures; our critique
targets specifically \emph{plain vector-similarity retrieval without
structured metadata}. Hybrid systems with metadata filtering,
re-ranking, or graph-based retrieval {[}9{]} address some limitations.
Ontology-grounded retrieval can be understood as a strict, typed,
auditable subclass of the broader RAG family --- not a rejection of the
paradigm.

\subsubsection{2.3 Security Ontologies and AppSec
Core}\label{security-ontologies-and-appsec-core}

Several ontologies have been proposed for security knowledge
representation. CWE {[}11{]} and CAPEC {[}12{]} model the \emph{problem
space} (what can go wrong); AppSec Core {[}7{]} models the
\emph{solution space} (what practitioners do to prevent it). This
distinction is relevant for code generation: a developer needs to know
what to \emph{do}, not what can go \emph{wrong}. OWASP's Secure Coding
Practices {[}13{]} provides an unstructured checklist lacking the typed,
weighted structure needed for reproducible retrieval and verification.

AppSec Core v0 {[}7{]} provides a normalization ontology for application
security: 10 domain slices, 234 typed instances (ControlObjective,
Practice, Mechanism, Artifact), with structural invariance across all
slices. The companion paper {[}7{]} demonstrates that heterogeneous
framework requirements (SSDF, ASVS, SLSA, CIS Controls, and CAPEC)
converge on shared objectives within this ontology.

The ontology is supported by a structural index of 4,139 units, each
annotated with \texttt{document\_role}, \texttt{normative\_weight}, and
\texttt{heading\_path}. This metadata is what makes ontology-grounded
retrieval possible --- and what plain RAG lacks.

\subsubsection{2.4 Model Context Protocol}\label{model-context-protocol}

MCP {[}8{]} provides a standardized interface for LLM agents to access
external tools and knowledge sources. In this work, MCP is the
\textbf{delivery mechanism}, not the contribution. The contribution is
the ontological structure of the knowledge being delivered and the
retrieval method that produces auditable context.

\begin{center}\rule{0.5\linewidth}{0.5pt}\end{center}

\subsection{3. The Problem: Six Missing Properties in Plain
Retrieval}\label{the-problem-six-missing-properties-in-plain-retrieval}

This section identifies six properties that security-critical code
generation requires and that plain vector-similarity RAG does not
provide. The critique is scoped to retrieval without structured
metadata.

\subsubsection{3.1 No First-Class
Provenance}\label{no-first-class-provenance}

Vector-similarity RAG retrieves text chunks ranked by similarity. The
LLM receives content but not the identity of the requirement it
represents. Post-hoc, there is no structured mechanism to determine
which specific requirement influenced the generated code. For security
compliance, where audit trails require traceable links between
requirements and implementations, this is insufficient.

\subsubsection{3.2 No Explicit Typing}\label{no-explicit-typing}

Plain RAG treats all retrieved chunks as equivalent: a normative
requirement, an operational practice, an implementation mechanism, and
an anti-pattern are delivered to the LLM as undifferentiated text. There
is no structural distinction between a mandatory constraint (``input
MUST be validated by allowlist'') and an informational example
(``consider using JSON Schema at the API gateway''). The LLM cannot
determine from the chunk alone whether to treat content as a
non-negotiable requirement or optional guidance.

\subsubsection{3.3 No Epistemic Weight}\label{no-epistemic-weight}

Related but distinct from typing, epistemic weight concerns the
\emph{authority level} of a unit, not its ontological category. Plain
RAG has no mechanism to signal that one chunk carries mandatory
normative force (strong weight) while another is operational guidance
(medium weight). Security knowledge is heterogeneous in this dimension:
a ControlObjective, a Practice, and a Mechanism may all be topically
relevant, but only the ControlObjective is non-negotiable. Without
explicit weighting, the LLM has no basis to prioritize accordingly.

\subsubsection{3.4 No Reproducible
Grounding}\label{no-reproducible-grounding}

The same prompt may retrieve different chunks depending on the embedding
model version, index state, or similarity threshold. The security
grounding of code generation is therefore non-reproducible: two runs may
produce code grounded in different requirement subsets. For functional
code this is acceptable; for security, where the question is ``were all
applicable requirements considered?'', it undermines the audit trail.

\subsubsection{3.5 No Completeness
Guarantee}\label{no-completeness-guarantee}

Vector similarity retrieves the \emph{most similar} chunks, not
\emph{all relevant} ones. If a knowledge base contains seven input
validation requirements (VAL-001→007), RAG may retrieve three or four,
missing the rest. There is no mechanism to verify that the full
requirement family was considered.

This retrieval incompleteness is structurally analogous to the
\emph{claim gap} phenomenon in compliance knowledge management {[}10{]}:
in both cases, a partial surface is mistaken for the applicable
knowledge base.

\subsubsection{3.6 No Verification Against
Grounding}\label{no-verification-against-grounding}

Plain RAG provides no way to verify generated code against the same
requirements that grounded it. The retrieved chunks are consumed by the
LLM and discarded; there is no structured record of what was provided,
and no mechanism to check whether the output satisfies it.

\begin{center}\rule{0.5\linewidth}{0.5pt}\end{center}

\subsection{4. Method: Ontology-Grounded
Retrieval}\label{method-ontology-grounded-retrieval}

\subsubsection{4.1 Overview}\label{overview}

Ontology-grounded retrieval replaces the embed-search-inject pipeline
with a five-stage structured query pipeline:

\begin{enumerate}
\def\labelenumi{\arabic{enumi}.}
\tightlist
\item
  \textbf{Intent parsing}: the prompt is analyzed to identify relevant
  ontology slices
\item
  \textbf{Structured retrieval}: the knowledge graph is queried with
  slice, risk level, and document role filters, returning typed units
  --- not similarity-ranked chunks
\item
  \textbf{Context assembly}: retrieved units are ordered by epistemic
  weight (strong requirements first, then practices, then mechanisms),
  each with an explicit provenance header
\item
  \textbf{Grounded generation}: the LLM generates code with instructions
  to cite which requirements it satisfies
\item
  \textbf{Verification}: the same structured index is queried to verify
  the output
\end{enumerate}

The method is general: it requires (a) a typed ontology with explicit
relations, (b) a structural index with role and weight annotations, and
(c) a delivery protocol. We instantiate it with AppSec Core {[}7{]} and
MCP {[}8{]}, but the method is not bound to these components.

\subsubsection{4.2 Six Distinguishing
Properties}\label{six-distinguishing-properties}

\textbf{Table 1.} Property comparison: plain vector-similarity RAG
vs.~ontology-grounded retrieval.

\begin{longtable}{@{}>{\RaggedRight\arraybackslash}p{2.55cm}>{\RaggedRight\arraybackslash}p{5.2cm}>{\RaggedRight\arraybackslash}p{5.2cm}@{}}
\toprule
\textbf{Property} & \textbf{Vector-similarity RAG} & \textbf{Ontology-grounded} \\
\midrule
\endfirsthead
\toprule
\textbf{Property} & \textbf{Vector-similarity RAG} & \textbf{Ontology-grounded} \\
\midrule
\endhead
\bottomrule
\endfoot
\bottomrule
\endlastfoot
\textbf{Provenance} & Implicit (chunk ID, opaque) & Explicit and first-class (\texttt{document\_role} + \texttt{heading\_path}) \\
\textbf{Typing} & Absent or weak; all chunks treated equivalently & Explicit: ControlObjective, Practice, Mechanism, Artifact (AppSec Core v0); AntiPattern, Requirement, Signal (runtime layer) \\
\textbf{Weighting} & Absent; no epistemic differentiation & First-class: \texttt{normative\_weight} (strong / medium / low) \\
\textbf{Reproducibility} & Non-reproducible (embedding-dependent) & Reproducible given fixed classification and index version \\
\textbf{Completeness} & Not assessable (similarity $\neq$ coverage) & Assessable via retrieval contract invariants \\
\textbf{Verifiability} & Not against grounding (no structured record) & Same structured index for generation and verification \\
\end{longtable}

These properties are not binary (RAG has none / ontology has all). They
represent a spectrum; Table 1 describes the \emph{baseline} forms of
each approach.

\subsubsection{4.3 Reproducible Grounding}\label{reproducible-grounding}

The retrieval path has two stages with different reproducibility
properties. First, the prompt is classified into a slice set; this stage
is probabilistic because it depends on the intent parser. Second, the
activated slice set is used to issue a structured query over the index;
given the same classification, ontology version, and index version, this
retrieval stage is reproducible.

We distinguish \textbf{classification uncertainty} (mapping a prompt to
ontology slices --- non-deterministic) from \textbf{retrieval
reproducibility} (querying the structural index with fixed filters ---
fully reproducible). Given the same classification, ontology version,
and index version, the same requirements are retrieved.

We use the term \textbf{reproducible grounding} rather than
``deterministic.'' The retrieval is reproducible \emph{conditioned on a
fixed classification}. The classification step introduces uncertainty
analyzed in Section 7.2.

\subsubsection{4.4 Typed Context Delivery}\label{typed-context-delivery}

Each unit delivered to the LLM carries its ontological type. The first
four types belong to the AppSec Core v0 ontology layer {[}7{]};
AntiPattern belongs to the runtime layer of the SbD-ToE knowledge graph,
which is projected through AppSec Core objectives:

\begin{itemize}
\tightlist
\item
  \textbf{ControlObjective} (strong): non-negotiable requirements
\item
  \textbf{Practice} (strong): operational approaches implementing
  objectives
\item
  \textbf{Mechanism} (medium): technical implementation patterns (LLM
  may choose alternatives)
\item
  \textbf{Artifact} (medium): evidence requirements (output should be
  testable)
\item
  \textbf{AntiPattern} (medium): explicit prohibitions (runtime layer)
\end{itemize}

The following structured context envelope illustrates what the LLM
receives for one requirement:

\begin{Shaded}
\begin{Highlighting}[]
\FunctionTok{slice}\KeywordTok{:}\AttributeTok{ ACO{-}IVF}
\FunctionTok{risk\_level}\KeywordTok{:}\AttributeTok{ L2}
\FunctionTok{requirement}\KeywordTok{:}
\AttributeTok{  }\FunctionTok{id}\KeywordTok{:}\AttributeTok{ VAL{-}003}
\AttributeTok{  }\FunctionTok{type}\KeywordTok{:}\AttributeTok{ ControlObjective}
\AttributeTok{  }\FunctionTok{weight}\KeywordTok{:}\AttributeTok{ strong}
\AttributeTok{  }\FunctionTok{description}\KeywordTok{:}\AttributeTok{ }\StringTok{"Schema validation before deserialization"}
\AttributeTok{  }\FunctionTok{validation\_method}\KeywordTok{:}\AttributeTok{ }\StringTok{"Verify typed model or JSON Schema at API entry point"}
\AttributeTok{  }\FunctionTok{expected\_evidence}\KeywordTok{:}\AttributeTok{ }\StringTok{"Test suite with malformed payload cases"}
\AttributeTok{  }\FunctionTok{provenance}\KeywordTok{:}
\AttributeTok{    }\FunctionTok{document\_role}\KeywordTok{:}\AttributeTok{ requirements\_catalog}
\AttributeTok{    }\FunctionTok{normative\_weight}\KeywordTok{:}\AttributeTok{ strong}
\AttributeTok{    }\FunctionTok{heading\_path}\KeywordTok{:}\AttributeTok{ }\StringTok{"Validação de Requisitos \textgreater{} VAL{-}003"}
\end{Highlighting}
\end{Shaded}

Two identifier namespaces appear in this envelope. \texttt{ACO-IVF} is
an AppSec Core v0 slice identifier; \texttt{ACO-IVF-*} would be the
native AppSec Core objective IDs. \texttt{VAL-003} is a corpus-level
requirement-family identifier from the SbD-ToE manual layer --- it is
the provenance ID of the requirement as it exists in the structural
index, mapped to the AppSec Core objective space through the projection
layer. The \texttt{heading\_path} field similarly carries corpus
provenance. This distinction matters for the verification loop:
completeness is assessed against AppSec Core ControlObjectives; corpus
IDs provide the traceable source reference.

This envelope enables the LLM to: (a) treat the requirement as
non-negotiable; (b) cite it in output (\texttt{//\ VAL-003}); (c) know
that evidence (test suite) should accompany the implementation. The same
envelope is available to the verification loop.

\subsubsection{4.5 The Retrieval Contract}\label{the-retrieval-contract}

We formalize the retrieval as a contract:

\begin{quote}
\textbf{Definition (Retrieval Contract).}

\emph{Input:} prompt \emph{p}, risk level \emph{l} ∈ \{L1, L2, L3\},
optional explicit slice hints \emph{H}

\emph{Output:} retrieved requirement set \emph{R} = \{(\emph{id},
\emph{type}, \emph{weight}, \emph{provenance})\}

\emph{Completeness invariant (conditional):} for every activated slice
\emph{s} and every ControlObjective \emph{co} ∈ \emph{s} where
\emph{co.risk\_level} ≤ \emph{l}: \emph{co} ∈ \emph{R}. All mandatory
requirements for activated slices at the requested risk level must be
returned. This invariant is conditional on correct slice activation;
completeness is guaranteed with respect to the activated ontology
region, not with respect to the user prompt in an unrestricted sense.

\emph{Provenance invariant:} every element of \emph{R} carries a
\texttt{(document\_role,\ normative\_weight,\ heading\_path)} triple
traceable to the structural index.
\end{quote}

The completeness invariant is testable: given a slice activation and
risk level, one can enumerate expected ControlObjectives and verify the
retrieval returned all of them. This is not possible with
similarity-ranked retrieval, where ``enough'' is undefined.

The completeness invariant is scoped to \textbf{ControlObjectives}
within activated slices. The richer payload delivered to the LLM ---
Practices, Mechanisms, AntiPatterns, and corpus-level identifiers --- is
auxiliary delivery around an objective-complete core. This auxiliary
content does not expand the completeness guarantee; it enriches the
grounding context for generation and verification without altering the
invariant's scope.

\subsubsection{4.6 Proportional
Application}\label{proportional-application}

Risk level is a retrieval parameter, not a post-hoc filter:

\begin{itemize}
\tightlist
\item
  \textbf{L1 (low risk):} L1-mandatory requirements only
\item
  \textbf{L2 (standard):} L1 + L2-mandatory requirements
\item
  \textbf{L3 (high risk):} all requirements including L3-specific and
  recommended practices
\end{itemize}

This prevents over-engineering low-risk components and ensures high-risk
components receive full treatment.

\begin{center}\rule{0.5\linewidth}{0.5pt}\end{center}

\subsection{5. Worked Examples}\label{worked-examples}

The following examples are \textbf{illustrative demonstrations of
traceability properties}, not comparative evaluations or standalone
reproducibility packages. Each traces the intended pipeline from prompt
to verification in compact form. Failure modes are discussed in Section
7.

For contrast, a plain vector-similarity RAG pipeline proceeds from
prompt to embedding, similarity-ranked chunks, and then generation,
without a structured record linking the output back to specific
requirements. The ontology-grounded pipeline instead proceeds from
prompt to slice classification, structured retrieval, typed requirement
delivery, generation, and verification against the same index, thereby
preserving a closed audit trail.

\subsubsection{5.1 API Input Validation}\label{api-input-validation}

\textbf{Prompt:} ``Write a Python API endpoint that accepts a JSON
payload with user profile data.''

\textbf{Retrieval:} ACO-IVF slice at L2. Returns (IDs of the form
\texttt{VAL-*}, \texttt{ERR-*} are corpus requirement-family identifiers
from the SbD-ToE runtime layer; \texttt{ACO-IVF} is the AppSec Core v0
slice):

\begin{longtable}{@{}L{3.1cm}L{1.35cm}L{1.55cm}L{6.95cm}@{}}
\toprule
\textbf{Type} & \textbf{Weight} & \textbf{ID} & \textbf{Requirement} \\
\midrule
\endfirsthead
\toprule
\textbf{Type} & \textbf{Weight} & \textbf{ID} & \textbf{Requirement} \\
\midrule
\endhead
\bottomrule
\endfoot
\bottomrule
\endlastfoot
ControlObjective & strong & VAL-001 & All input validated by allowlist \\
ControlObjective & strong & VAL-003 & Schema validation before deserialization \\
ControlObjective & strong & VAL-005 & Validation before internal use \\
ControlObjective & strong & ERR-001 & Fail-safe; reject invalid input without revealing internals \\
Practice & medium & --- & Use typed model at API entry point \\
Mechanism & medium & --- & \texttt{pydantic BaseModel} with \texttt{strict=True} \\
AntiPattern & medium & --- & Never use \texttt{eval()} / \texttt{exec()} on user input \\
\end{longtable}

\textbf{Generated code} (with requirement citations):

\begin{Shaded}
\begin{Highlighting}[]
\ImportTok{from}\NormalTok{ pydantic }\ImportTok{import}\NormalTok{ BaseModel, ConfigDict}

\KeywordTok{class}\NormalTok{ UserProfile(BaseModel):}
\NormalTok{    model\_config }\OperatorTok{=}\NormalTok{ ConfigDict(strict}\OperatorTok{=}\VariableTok{True}\NormalTok{)  }\CommentTok{\# VAL{-}003: schema before deserialization}
\NormalTok{    name: }\BuiltInTok{str}
\NormalTok{    email: }\BuiltInTok{str}
\NormalTok{    age: }\BuiltInTok{int}

\AttributeTok{@app.post}\NormalTok{(}\StringTok{"/profile"}\NormalTok{)}
\ControlFlowTok{async} \KeywordTok{def}\NormalTok{ create\_profile(profile: UserProfile):}
    \CommentTok{\# VAL{-}001: allowlist validation via typed model fields}
    \CommentTok{\# VAL{-}005: validated before any internal use}
    \CommentTok{\# ERR{-}001: pydantic returns 422 without exposing internals}
    \ControlFlowTok{return} \ControlFlowTok{await}\NormalTok{ save\_profile(profile)}
\end{Highlighting}
\end{Shaded}

\textbf{Verification:} VAL-001 ✓, VAL-003 ✓, VAL-005 ✓, ERR-001 ✓,
AntiPattern (no \texttt{eval}) ✓. All retrieved requirements addressed.
Audit trail complete.

\subsubsection{5.2 Kubernetes Container
Hardening}\label{kubernetes-container-hardening}

\textbf{Prompt:} ``Write a Kubernetes deployment manifest for a web
application.''

\textbf{Retrieval:} ACO-SPC + ACO-RPR at L2. Returns CFG-001 (no debug),
CFG-003 (no hardcoded credentials), securityContext practices, and
\texttt{privileged:\ true} anti-pattern.

\textbf{Generated manifest} includes: \texttt{runAsNonRoot:\ true},
\texttt{allowPrivilegeEscalation:\ false},
\texttt{readOnlyRootFilesystem:\ true}, credentials from
\texttt{secretKeyRef}, no debug ports exposed. Each security property
traceable to a retrieved requirement.

\subsubsection{5.3 CI/CD Pipeline
Security}\label{cicd-pipeline-security}

\textbf{Prompt:} ``Write a GitHub Actions workflow for building and
deploying a container image.''

\textbf{Retrieval:} ACO-SCBI + ACO-RPR at L2. Returns provenance
attestation, artifact signing, runner isolation, secret injection, and
branch protection requirements.

\textbf{Generated workflow} includes: attestation generation, cosign
signing, \texttt{secrets.*} only, image reference by digest, and
approval gate for production. Each step traceable to retrieved
requirements.

\begin{center}\rule{0.5\linewidth}{0.5pt}\end{center}

\subsection{6. The Verification Loop}\label{the-verification-loop}

The verification loop closes the audit trail. We organize verification
into three classes by automation level.

\subsubsection{6.1 Syntactic Verification (Fully
Automatable)}\label{syntactic-verification-fully-automatable}

Checks on presence or absence of concrete constructs: anti-pattern
absence (\texttt{eval()}, \texttt{privileged:\ true}), configuration
presence (\texttt{runAsNonRoot:\ true}), citation validity (does
\texttt{//\ VAL-003} exist in the knowledge graph and in the retrieved
set?). These are deterministic.

\subsubsection{6.2 Semantic Approximation
(Semi-Automatable)}\label{semantic-approximation-semi-automatable}

Checks requiring limited static analysis: data flow (does user input
pass through validation before a database query?), schema enforcement
(is strict mode enabled?), citation plausibility (does code near
\texttt{//\ VAL-003} actually perform schema validation?). Approximable
with SAST tools configured from the same requirement families.

\subsubsection{6.3 Human Judgment
Required}\label{human-judgment-required}

Checks requiring human review: architectural adequacy (least
privilege?), business logic constraints (appropriate fail-safe?),
defense completeness (all attack surfaces?). The structured audit trail
makes human review tractable but not automatable.

\subsubsection{6.4 The Closed Loop}\label{the-closed-loop}

Generation and verification use the \textbf{same structured index}:

\begin{itemize}
\tightlist
\item
  No requirements lost in translation between grounding and verification
\item
  Verification scoped to exactly the relevant requirements
\item
  Audit trail complete from prompt, through slice selection and
  retrieved requirements, to generated code and verification results
\end{itemize}

This closed loop is not achievable with plain RAG, where retrieved
chunks are consumed without structured record.

\begin{center}\rule{0.5\linewidth}{0.5pt}\end{center}

\subsection{7. Discussion}\label{discussion}

\subsubsection{7.1 What is Reproducible, What is
Not}\label{what-is-reproducible-what-is-not}

Ontology-grounded retrieval makes the \emph{grounding} reproducible, not
the \emph{output}. Once slice activation is fixed, the same query yields
the same requirements. The LLM's generation remains probabilistic. But
the security contract is fixed: same requirements, same verification
criteria, same audit trail.

\subsubsection{7.2 The Intent Parser: Primary Source of
Uncertainty}\label{the-intent-parser-primary-source-of-uncertainty}

The intent parser --- mapping a prompt to ontology slices --- is the
primary source of uncertainty. If it assigns wrong slices, retrieval is
reproducible but \emph{reproducibly wrong}: the completeness invariant
holds for activated slices, but the slices themselves may be incorrect.

Mitigation strategies (in order of robustness):

\begin{enumerate}
\def\labelenumi{\arabic{enumi}.}
\tightlist
\item
  \textbf{Explicit tagging} (\texttt{@ACO-IVF}) --- bypasses the parser;
  most reliable, least ergonomic
\item
  \textbf{Conservative over-retrieval} (top-3 slices) --- trades context
  space for coverage safety
\item
  \textbf{Confidence thresholding} --- low-confidence slices trigger
  fallback to tagging or over-retrieval
\end{enumerate}

Empirical characterization of parser accuracy is future work.

\subsubsection{7.3 Cost and
Generalizability}\label{cost-and-generalizability}

Ontology-grounded retrieval requires a maintained, typed knowledge graph
--- a non-trivial investment compared to indexing unstructured
documentation for RAG. The cost is justified when security verification
properties are required. For non-security contexts, standard RAG may
suffice.

The method is general in structure: it requires a typed ontology, a
stratified index, and a delivery protocol. We instantiate it for AppSec
with AppSec Core and MCP. The pattern is in principle applicable to
other domains where structured, verifiable grounding is needed ---
regulatory compliance, medical guidelines, safety-critical engineering
--- but this broader applicability remains a methodological claim, not a
demonstrated result.

\subsubsection{7.4 Relationship to Static
Analysis}\label{relationship-to-static-analysis}

Ontology-grounded retrieval is complementary to SAST:

\begin{itemize}
\tightlist
\item
  \textbf{Requirement families → SAST rule selection}: activating a
  slice enables the corresponding rule sets
\item
  \textbf{AntiPatterns → lint rules}: each prohibition translates to a
  static analysis check
\item
  \textbf{Validation methods → test templates}: each requirement's
  \texttt{validation\_method} seeds test generation
\end{itemize}

Both generation and verification operate from the same structured
requirements, ensuring consistency.

\subsubsection{7.5 Illustrative Property
Snapshot}\label{illustrative-property-snapshot}

We report a small illustrative comparison of retrieval properties for
the three prompts in Section 5. This is not a controlled evaluation and
should not be read as evidence of superiority in code quality or
security outcomes; it serves only to make the contract properties
operationally concrete. The baseline shown here is a single illustrative
instantiation using one embedding model over the same corpus without
structured metadata:

\textbf{Table 2.} Illustrative retrieval property snapshot for worked
examples (Section 5). Not a controlled evaluation --- see §8.

\begin{longtable}{@{}>{\RaggedRight\arraybackslash}p{4.5cm}>{\RaggedRight\arraybackslash}p{4.1cm}>{\RaggedRight\arraybackslash}p{4.6cm}@{}}
\toprule
\textbf{Metric} & \textbf{Vector-similarity RAG} & \textbf{Ontology-grounded} \\
\midrule
\endfirsthead
\toprule
\textbf{Metric} & \textbf{Vector-similarity RAG} & \textbf{Ontology-grounded} \\
\midrule
\endhead
\bottomrule
\endfoot
\bottomrule
\endlastfoot
Requirements retrieved (of applicable¹) & 3--5 of 7 (varies) & 7 of 7 (stable) \\
Requirements with explicit ID in context & 0 & 7 \\
Epistemic weight attached & No & Yes (strong/medium) \\
Verification checks executable post-gen & 0 (no record) & 5--7 (syntactic + semantic) \\
Audit trail constructible & No & Yes \\
\end{longtable}

¹ ``Applicable'' = all ControlObjectives in activated slices with
\texttt{risk\_level} ≤ selected level (per retrieval contract, Section
4.5).

This snapshot compares \emph{retrieval properties} of a single baseline
instantiation, not code quality or security outcomes. A rigorous
evaluation with multiple models, prompts, annotators, and statistical
controls is future work.

\subsubsection{7.6 What This is Not}\label{what-this-is-not}

This method does not claim that LLMs should write security-critical code
autonomously. Human review remains essential. What ontology-grounded
retrieval provides is \textbf{structure for that review}: reviewers see
exactly which requirements grounded the generation and can verify each
one.

Nor does it guarantee secure code. The claim is narrower:
ontology-grounded retrieval produces \textbf{auditable,
requirement-traceable grounding}. The security of the output depends on
the requirements' quality, the LLM's adherence, and the verification
loop's thoroughness.

\begin{center}\rule{0.5\linewidth}{0.5pt}\end{center}

\subsection{8. Limitations}\label{limitations}

\textbf{No controlled experiment.} The validation is methodological
(compact illustrative scenarios making the contract properties
concrete), not experimental. A comparative study measuring vulnerability
rates and requirement coverage across grounding modes is future work.

\textbf{Intent parser uncharacterized.} The primary source of
uncertainty has no accuracy data. Future work should measure
classification precision and recall per slice across prompt types.

\textbf{Single instantiation.} The method is demonstrated with AppSec
Core and MCP. Generalization to other ontologies and delivery mechanisms
is argued but not validated.

\textbf{Worked examples are happy paths.} Failure modes (wrong slice,
incomplete retrieval, LLM non-compliance) are discussed but not
systematically characterized.

\begin{center}\rule{0.5\linewidth}{0.5pt}\end{center}

\subsection{9. Conclusion}\label{conclusion}

We have presented ontology-grounded retrieval as a method for auditable
LLM-assisted secure code generation. The method replaces probabilistic
vector similarity with reproducible structured queries against a typed
security knowledge graph, providing six properties that plain RAG does
not inherently offer: first-class provenance, explicit typing, epistemic
weighting, reproducible grounding, completeness assessment, and
verifiability against the same index.

The retrieval contract formalizes completeness and provenance invariants
that make grounding testable. The three-class verification taxonomy
makes automation boundaries explicit. The closed verification loop ---
where the same index grounds generation and verifies output --- closes
the audit trail.

The method does not make LLM-generated code secure. It makes the
grounding auditable: requirements are explicit, provenance is traceable,
verification is reproducible. The method is general in structure but
validated here only in an AppSec-specific instantiation; its broader
applicability remains a methodological claim to be tested empirically in
future work.

In a landscape where LLMs generate increasing volumes of code, the
ability to say ``this code was grounded in these specific requirements,
and here is the evidence'' is a meaningful step toward responsible
adoption.

\begin{center}\rule{0.5\linewidth}{0.5pt}\end{center}

\subsection{10. Artifact Availability}\label{artifact-availability}

Curated supporting artifacts for this paper are available in the
companion public repository at
\url{https://github.com/sbd-ai-runtime/appsec-core-ontology-research}.
For this paper, the relevant materials are organized under
\texttt{papers/03-ontology-grounded-retrieval/artifacts/}, notably
\texttt{papers/03-ontology-grounded-retrieval/artifacts/retrieval\_contract/}
and
\texttt{papers/03-ontology-grounded-retrieval/artifacts/runtime\_snapshot/},
which provide the released retrieval contract and runtime-grounding
snapshot referenced by the method. The same repository also contains
this paper's curated source under
\texttt{papers/03-ontology-grounded-retrieval/source/} and its public
PDF under \texttt{papers/03-ontology-grounded-retrieval/pdf/}.

\begin{center}\rule{0.5\linewidth}{0.5pt}\end{center}

\subsection{References}\label{references}

{[}1{]} H. Pearce, B. Ahmad, B. Tan, B. Dolan-Gavitt, and R. Karri,
``Asleep at the keyboard? Assessing the security of GitHub Copilot's
code contributions,'' in \emph{Proc. 43rd IEEE S\&P}, 2022,
pp.~754--768, doi: 10.1109/SP46214.2022.9833571.

{[}2{]} N. Tihanyi, T. Bisztray, M. A. Ferrag, R. Jain, and L. C.
Cordeiro, ``How secure is AI-generated code: A large-scale comparison of
large language models,'' \emph{Empirical Softw. Eng.}, vol.~30, art. 47,
2025, doi: 10.1007/s10664-024-10590-1.

{[}3{]} C. Tony, M. Mutas, N. Díaz Ferreyra, and R. Scandariato,
``LLMSecEval: A dataset for security evaluation of LLM code
suggestions,'' in \emph{Proc. MSR}, 2023, pp.~588--592, doi:
10.1109/MSR59073.2023.00084.

{[}4{]} C. Tony, N. E. Díaz Ferreyra, M. Mutas, S. Dhif, and R.
Scandariato, ``Prompting techniques for secure code generation: A
systematic investigation,'' \emph{ACM Trans. Softw. Eng. Methodol.},
vol.~34, no. 8, art. 225, 2025, doi: 10.1145/3722108.

{[}5{]} M. Chen et al., ``Evaluating large language models trained on
code,'' arXiv:2107.03374, 2021.

{[}6{]} P. Lewis et al., ``Retrieval-augmented generation for
knowledge-intensive NLP tasks,'' in \emph{NeurIPS}, vol.~33, 2020,
pp.~9459--9474.

{[}7{]} P. Farinha, ``AppSec Core: A normalized ontology for security
requirements across heterogeneous frameworks,'' preprint, 2026. {[}arXiv
ID: TBD --- submitted as companion preprint{]}

{[}8{]} Anthropic, ``Model Context Protocol specification,'' version
2024-11-05, 2024. GitHub:
https://github.com/modelcontextprotocol/modelcontextprotocol/releases
(accessed: 2026-04-06).

{[}9{]} B. Peng et al., ``Graph retrieval-augmented generation: A
survey,'' \emph{ACM Trans. Inf. Syst.}, 2025, doi: 10.1145/3777378.
{[}DOI to be verified against published record{]}

{[}10{]} P. Farinha, ``Coverage-preserving compilation of normative and
empirical security knowledge,'' preprint, 2026. {[}arXiv ID: TBD ---
submitted as companion preprint{]}

{[}11{]} MITRE. Common Weakness Enumeration (CWE). Available:
https://cwe.mitre.org

{[}12{]} MITRE. Common Attack Pattern Enumeration and Classification
(CAPEC). Available: https://capec.mitre.org

{[}13{]} OWASP. Secure Coding Practices Quick Reference Guide, 2010.

\end{document}
